% !TeX root = ../教材题库_全部条目_转写工作稿.tex
% Source: XJ-B1-C05-RV-P206-Y1539
% 第5章 / 章末复习 复习题五 / 教材 p.206、207 / PDF 第213、214页
% OCR未校对：公式与文字均可能有误；未完成数学TeX转写；禁止直接发给学生
% 原题及相邻原解截图（核对用）：../题目原图/XJ-B1-C05-RV-P206-Y1539-01.webp, ../题目原图/XJ-B1-C05-RV-P206-Y1539-02.webp
% 本文件只在转写完成后作为题干使用；OCR 条目尚不是可靠题干。
\begin{bankitem}{复习题五 第9题}{第5章 / 章末复习 复习题五 / 教材 p.206、207 / PDF 第213、214页}{XJ-B1-C05-RV-P206-Y1539}{OCR未校对：公式与文字均可能有误；未完成数学TeX转写；禁止直接发给学生}
\begin{ocrdraft}
[教材p.206 / PDF第213页]
\par
9．
\par
已知函数3，一2cos
\par
一3，
\par
[教材p.207 / PDF第214页]
\par
求函数的单调区间及取得最大、最小值时自变量x的集合；
\par
（2）判断函数的奇偶性．
\par
利用三角函数线求函数丆）==cos，在〔
\end{ocrdraft}
\end{bankitem}
